\documentclass[10pt,twocolumn,letterpaper]{article}
\usepackage[rebuttal]{cvpr}

% Include other packages here, before hyperref.
\usepackage{graphicx}
\usepackage{amsmath}
\usepackage{amssymb}
\usepackage{booktabs}
\usepackage{url}
\usepackage{microtype}

\usepackage[pagebackref,breaklinks,colorlinks,bookmarks=false]{hyperref}

% Support for easy cross-referencing
\usepackage[capitalize]{cleveref}
\crefname{section}{Sec.}{Secs.}
\Crefname{section}{Section}{Sections}
\Crefname{table}{Table}{Tables}
\crefname{table}{Tab.}{Tabs.}

%%%%%%%%% PAPER ID  - PLEASE UPDATE
\def\cvprPaperID{2006.09882} 
\def\confName{MAI}
\def\confYear{Fall 2026}

\begin{document}

%%%%%%%%% TITLE - PLEASE UPDATE
\title{Unsupervised Learning of Visual Features by Contrasting Cluster Assignments\\ B12508026 Tai, Wei-Hsuan}  % **** Enter the paper title here

\maketitle
\thispagestyle{empty}
\appendix

%%%%%%%%% BODY TEXT - ENTER YOUR RESPONSE BELOW

\section{Summary}
This paper proposes SwAV (Swapping Assignments between Views), an online clustering-based self-supervised method for visual representation learning.
Instead of comparing features of different views directly as in contrastive learning, it assigns each view to learnable prototypes and predicts the cluster assignment of one view from the feature of another view of the same image.
Together with multi-crop, a data augmentation strategy that adds low-resolution crops at little extra cost, SwAV achieves 75.3\% top-1 accuracy on ImageNet linear evaluation with ResNet-50 and surpasses supervised pretraining on all considered transfer tasks.

\section{Strengths}
\begin{enumerate}
    \item Memory efficient and flexible in batch size. Without a memory bank or momentum encoder, SwAV stores only 3{,}840 features (vs.\ 65{,}536 for MoCov2) and reaches 72.0\% with a batch size of 256 on 4 GPUs (Table~3).
    \item Effective multi-crop augmentation. Low-resolution crops add views at little extra cost and improve SimCLR, DeepCluster-v2, and SeLa-v2 by 2.4--4.6\% top-1 accuracy (Fig.~3).
    \item Truly online method. Codes are computed within each mini-batch, which enables pretraining on 1B uncurated Instagram images (Sec.~5.4).
    \item Comprehensive experiments. SwAV surpasses supervised pretraining on all five transfer tasks, with a best-effort fair comparison with SimCLR under identical settings (Fig.~3).
\end{enumerate}

\section{Major Weaknesses}
\begin{enumerate}
    \item The benefit of swapping is not isolated, and the online advantage is not quantified. Under identical settings (Fig.~3), SwAV is on par with DeepCluster-v2 (70.1 vs.\ 70.2; 75.3 vs.\ 75.2 in the best setting), which the paper itself views as a special case of swapping. No variant without swapping (e.g., predicting a view's own code) is reported, and the online advantage over DeepCluster-v2 is argued only by construction, without any cost or large-scale comparison.

    \item $K$, $\varepsilon$, and equipartition are not analyzed. $K$ ($\sim$3K) is fixed without justification, although it also sets the queue size and thus the claimed memory advantage. The value of $\varepsilon$ is not reported, although a large $\varepsilon$ leads to collapse (Sec.~3.1), and none of these is ablated in the main text. Equipartition also assumes balanced concepts, yet Sec.~5.4 gives no analysis of prototype usage on long-tailed uncurated data.

    \item (2026) Limited novelty and outdated baselines. Multi-crop is now standard in DINO~\cite{caron2021dino}, iBOT~\cite{zhou2022ibot}, and DINOv2~\cite{oquab2024dinov2}, and DINOv2 adopts Sinkhorn-Knopp centering. The paper compares only with 2019--2020 methods, not with these, MAE~\cite{he2022mae}, or CLIP~\cite{radford2021clip}. DINO, iBOT, DINOv2, and CLIP exceed 80\% top-1 with frozen ViT features, whereas DINO with the same ResNet-50 also reaches 75.3\%, so the gap lies in architecture and scale, which the paper does not explore.

    \item (2026) Image-level objective, CNN-only, and narrow evaluation. There is no patch-level objective as in iBOT or DINOv2, and all results use ResNets rather than ViTs. Frozen-feature segmentation, depth, retrieval, and robustness benchmarks, now standard for general-purpose visual features, are also missing.
\end{enumerate}

\section{Minor Weaknesses}
\begin{enumerate}
    \item Inconsistent compute claims. The abstract claims multi-crop adds no compute, while Sec.~3.2 admits ``only a small increase''. The 2$\times$224+6$\times$96 setting processes $\sim$1.55$\times$ the pixels of 2$\times$224, and no equal-budget comparison is given for the headline result.

    \item Clarity and typos. Below Eq.~(4), ``on average at least $B/K$ times'' should be exactly $B/K$ of the assignment mass, and the stop-gradient on $\mathbf{Q}$ is not stated. $\mathbf{z}_i$ before Eq.~(2) should be $\mathbf{z}_t$, and $\mathbf{c}_k$ should be $\mathbf{c}_K$ (Sec.~3.1). Published works (e.g., SimCLR, MoCo, PIRL) are cited as arXiv preprints. Typos: ``backpropragation'' (Fig.~1), ``uniformely'' and ``an unique'' (Sec.~3.1), ``this is line with'' (Sec.~4.2), ``correspondance'' and ``implementating'' (Sec.~5.1), ``accucary'' (Fig.~3).

    \item Missing settings. The value of $\varepsilon$, the crop scales, and when the queue starts to be used are not given, and the claim that 3 Sinkhorn iterations are ``sufficient'' is not supported by numbers.

    \item No variance over seeds. Given the 0.1--0.2\% gaps with DeepCluster-v2 and the $\sim$1.3-point detection margins over supervised pretraining, mean and standard deviation over a few seeds are needed.

    \item (2026) Reproducibility and domain shift. The 1B Instagram dataset is neither released nor documented, and total compute (about 3{,}200 V100 GPU-hours for the 800-epoch model) is not reported. The Broader Impact section cites medical imaging, but no such experiment is given; under equipartition (Major~2), data dominated by ``normal'' images is where the method is most uncertain.
\end{enumerate}

\section{Justification}
Although SwAV is only on par with DeepCluster-v2, it outperforms contrastive methods under identical settings, works online, and remains effective with small batches (Table~3). Major~1 and~2 concern analysis rather than the core results, so I suggest Weak Accept. A no-swap ablation and an analysis of $K$, $\varepsilon$, and equipartition in the rebuttal could raise my rating.

(2026) Its main ideas are now standard components of later methods (Major~3), and its ResNet-50 results are far below current ViT-based models. Since the lack of novelty cannot be addressed by additional experiments, I suggest Weak Reject from today's point of view.

\section{Final Rating}
4: Weak Accept

(2026) 2: Weak Reject

\section{Reflections}
At first I was skeptical of this paper: the results seemed to come from single runs, the benchmarks felt narrow, and the gain over DeepCluster-v2 was negligible. Looking into the 2020 context, I realized I was judging it by today's standards. One 800-epoch run already costs about 3{,}200 GPU-hours, and SwAV clearly beats contrastive methods under the same settings.

My main lesson came from seeing multi-crop adopted by DINO and Sinkhorn-Knopp centering by DINOv2. A paper's value lies less in its headline number than in whether its ideas survive in later methods. ResNet~\cite{he2016resnet} is a clear example, its ImageNet accuracy has long been surpassed, yet residual connections remain in almost every modern architecture, including Transformers~\cite{vaswani2017attention}.

For my own research, medical imaging tasks often suffer from limited data. I once doubted whether it is worth training a model for a disease too rare to have enough images. Yet SwAV, built for no particular downstream task, learned features from unlabeled images that transfer to tasks its authors never targeted. What a model learns can matter more than the task it was designed for.

%%%%%%%%% REFERENCES
{\small
\bibliographystyle{ieee_fullname}
\bibliography{egbib}


}

\end{document}
